% part: intuitionistic-logic
% chapter: propositions-as-types
% section: proofs-to-terms

\documentclass[../../../include/open-logic-section]{subfiles}

\begin{document}

\olfileid[hi]{int}{pty}{pt}

\olsection{\usetoken{P}{derivation} को प्रमाण पदों में बदलना}

प्राकृतिक निगमन के !!{proof}s को प्रमाण पदों में बदलने की प्रक्रिया, अर्थात्
~\Log{N2i} के !!{proof}s का~\Log{tN2i} के !!{proof}s में अनुवाद, अत्यंत सीधा
है। इसमें केवल !!{derivation} के किसी अनुक्रम की दाईं ओर प्रत्येक !!{formula}
के बाईं ओर प्रमाण पद लिखना है, जिससे $\Gamma \Sequent M : !A$ रूप के अनुक्रम
मिलते हैं। तब हम कहेंगे कि प्रसंग~$\Gamma$ में $M$,~$!A$ का \emph{साक्षी} है।
कौन-सा प्रमाण पद लिखना है, यह पूर्णतः \Log{tN2i} के नियम निर्धारित करते हैं।

\begin{prop}
\Log{N2i} में $\Gamma \Sequent !A$ का प्रत्येक !!{proof}~$\delta$,
\Log{tN2i} में $\Gamma \Sequent N : !A$ के किसी !!a{proof}~$\delta'$ के अनुरूप
है। प्रमाण पद $N$,~$\delta$ से अद्वितीय रूप से निर्धारित होता है।
\end{prop}

\begin{proof}
हम~$\delta$ की ऊँचाई पर आगमन से आगे बढ़ते हैं। यदि $\delta$ अभिगृहीत $x:!A
\Sequent !A$ है, तो $\delta'$ अभिगृहीत~$x:!A \Sequent x:!A$ है। यदि $\delta$
किसी अनुमान-नियम पर समाप्त होता है, तो आगमन परिकल्पना प्रत्येक आधार-वाक्य का
प्रमाण पद और प्रत्येक आधार-वाक्य के अनुरूप, उन प्रमाण पदों वाले अनुक्रमों के
\Log{tN2i}-!!{proof}s देती है। हम \Log{tN2i} का संगत नियम लगाते हैं; संगत
प्रमाण पद \Log{tN2i} का अनुमान-नियम निर्धारित करता है।
\end{proof}

\begin{ex}
  $(!A \land !B) \lif (!B \land !A)$ के \Log{N1i} में इस अत्यंत सरल
  !!{proof} पर विचार करें:
  \[
  \AxiomC{$\Discharge{!A\land !B}{x}$}
  \RightLabel{\Elim\land[2]}
  \UnaryInfC{$!B$}
  \AxiomC{$\Discharge{!A\land !B}{x}$}
  \RightLabel{\Elim\land[1]}
  \UnaryInfC{$!A$}
  \RightLabel{\Intro\land}
  \BinaryInfC{$!B \land !A$}
  \DischargeRule{\Intro\lif}{x}
  \UnaryInfC{$(!A \land !B) \lif (!B \land !A)$}
  \DisplayProof
  \]
  \Log{N2i} में इसका रूप यह है:
  \[
  \Axiom$x : !A\land !B \fCenter !A\land !B$
  \RightLabel{\Elim\land[2]}
  \UnaryInf$x : !A\land !B \fCenter !B$
  \Axiom$x : !A\land !B \fCenter !A\land !B$
  \RightLabel{\Elim\land[1]}
  \UnaryInf$x : !A\land !B \fCenter !A$
  \RightLabel{\Intro\land}
  \BinaryInf$x : !A\land !B \fCenter !B \land !A$
  \RightLabel{\Intro\lif}
  \UnaryInf$\fCenter (!A \land !B) \lif (!B \land !A)$
  \DisplayProof
  \]
  हम ऊपर से नीचे की ओर प्रमाण पद देते हैं। अभिगृहीतों के प्रमाण पद प्रसंग के
  चर~$x$ ही हैं। $\Elim\land[i]$ से संबंधित प्रमाण पद क्रमशः $\proj{2}{x}$
  और~$\proj{1}{x}$ निर्धारित होते हैं। \Intro{\land} लगाने पर इन दोनों प्रमाण
  पदों को मिलाते हैं: $\pair{\proj{2}{x}}{\proj{1}{x}}$। अंत में
  \Intro{\lif} नियम कोई $\lambd$-अमूर्तन लगाता है,~$x$ को आबद्ध करता है और
  दर्ज करता है कि $x$ ने मान्यता~$!A \land !B$ को चिह्नित किया था:
  $\lambd[\typeof{x}{!A \land !B}][\pair{\proj{2}{x}}{\proj{1}{x}}]$। तब
  \Log{tN2i} में !!{proof} है:
  \[
  \Axiom$x : !A\land !B \fCenter x: !A\land !B$
  \RightLabel{\Elim\land[2]}
  \UnaryInf$x : !A\land !B \fCenter \proj{2}{x} : !B$
  \Axiom$x : !A\land !B \fCenter x: !A\land !B$
  \RightLabel{\Elim\land[1]}
  \UnaryInf$x : !A\land !B \fCenter \proj{1}{x} : !A$
  \RightLabel{\Intro\land}
  \BinaryInf$x : !A\land !B \fCenter \pair{\proj{2}{x}}{\proj{1}{x}} : 
    !B \land !A$
  \RightLabel{\Intro\lif}
  \UnaryInf$\fCenter \lambd[\typeof{x}{!A \land
  !B}][\pair{\proj{2}{x}}{\proj{1}{x}}] : (!A \land !B) \lif (!B \land !A)$
  \DisplayProof
  \]
\end{ex}


\end{document}
